% ============================================================================
% Section 5.3: Practical Deployment Recommendations
% ============================================================================
% Updated to reflect honest three-regime framing from crossover experiment.

\subsection{Practical Deployment Recommendations}
\label{sec:practical}

The prior degradation sweep (\S\ref{sec:metalearning_cost}) establishes
a crossover between strategies as prior quality varies.
This section translates those findings into deployment guidance,
including the \texttt{initial\_weights} parameter that lets
practitioners tune the trade-off between best-case and worst-case
performance.

% ----------------------------------------------------------------------------
% Recommendation 1: Strategy Selection
% ----------------------------------------------------------------------------
\paragraph{Strategy Selection.}
Figure~\ref{fig:prior_degradation}A provides the primary decision tool.
No single strategy dominates everywhere---each is optimal in a distinct
regime:

\begin{enumerate}[nosep]
    \item \textbf{Use Warmup-Only if priors are validated:}
    \begin{itemize}[nosep]
        \item Prior accuracy ${>}80\%$ on $N{=}100$--$200$ deployment
              samples
        \item Deployment distribution closely matches training data
              (PSI $<$0.20)
        \item \emph{Gain: up to 47\% lower regret than Corralling
              (30.0 vs 44.1 at $\alpha{=}0$)}
    \end{itemize}

    \item \textbf{Use Corralling if prior quality is unknown:}
    \begin{itemize}[nosep]
        \item No validation data available yet
        \item Deployment involves cross-domain or evolving distributions
        \item Model reliability may degrade (\S\ref{sec:catastrophic})
        \item Set initial warmup weight $w_0{=}0.7$ to reduce
              overhead by 11\% when priors are correct
        \item \emph{Gain: bounded worst-case regret (40--67 range vs
              warmup-only's 27--99)}
    \end{itemize}

    \item \textbf{Use Tabula Rasa if priors are known to be bad:}
    \begin{itemize}[nosep]
        \item Prior accuracy validated as ${<}50\%$ (adversarial)
        \item True cold start (no priors available at all)
        \item Severe domain mismatch confirmed
        \item \emph{Gain: 27\% lower regret than Corralling when
              warmup expert is harmful (49.2 vs 67.3 at $\alpha{=}1$)}
    \end{itemize}
\end{enumerate}

% ----------------------------------------------------------------------------
% Tuning the Trade-Off: Initial Weight Bias
% ----------------------------------------------------------------------------
\paragraph{Tuning the Trade-Off: \texttt{initial\_weights}.}
Corralling's overhead at $\alpha{=}0$ is not fixed---practitioners
can control it via the \texttt{initial\_weights} parameter.
This sets the meta-learner's starting allocation between the warmup
and tabula rasa experts, encoding \emph{prior trust}:

\begin{itemize}[nosep]
    \item $w_0{=}0.5$ (uniform): Maximum safety, no prior trust.
          Best-case regret 50; worst-case 60.
    \item $w_0{=}0.7$ (recommended): Moderate prior trust.
          Best-case regret 44 ($-$11\%); worst-case 67 ($+$9\%).
    \item $w_0{=}0.9$ (aggressive): Strong prior trust.
          Best-case regret 37 ($-$26\%); worst-case 89 ($+$49\%).
\end{itemize}

\noindent
The trade-off is \emph{asymmetric}: beyond $w_0{=}0.7$, the
worst-case cost grows superlinearly while the best-case improvement
saturates.
We recommend $w_0{=}0.7$ as the default; increase to 0.8--0.9 only
with high confidence in prior quality.

\begin{verbatim}
router = CorrallingRouter(
    experts=[warmup, tabula_rasa],
    initial_weights=np.array([0.7, 0.3]),  # prior trust
)
\end{verbatim}

% ----------------------------------------------------------------------------
% Why not always use Tabula Rasa?
% ----------------------------------------------------------------------------
\paragraph{Why Not Always Use Tabula Rasa?}
A natural question: if tabula rasa is prior-independent, why not always
use it?  Because validated priors are worth 39\% lower regret
(warmup-only's 30.0 vs tabula rasa's 49.2).
Tabula rasa sacrifices the \emph{upside} of good priors for
\emph{immunity} to bad ones.  It is the correct choice only when the
practitioner has positive evidence that priors are harmful.

% ----------------------------------------------------------------------------
% Adaptive deployment pattern
% ----------------------------------------------------------------------------
\paragraph{Adaptive Deployment Pattern.}
The library's \texttt{selection\_token} API supports real-time
strategy switching based on weight monitoring:

\begin{itemize}[nosep]
    \item Warmup weight ${>}0.8$ after 200--400 queries
          $\Rightarrow$ priors are strong; switch to warmup-only
    \item Balanced weights ($0.3$--$0.7$) $\Rightarrow$ genuine
          uncertainty; Corralling is appropriate
    \item Warmup weight ${<}0.3$ $\Rightarrow$ priors are harmful;
          switch to tabula rasa
    \item Weights frozen (variance ${<}0.01$) $\Rightarrow$
          implementation bug (missing \texttt{selection\_token})
\end{itemize}

This enables a \emph{start safe, optimise later} workflow: deploy with
Corralling, then transition to warmup-only or tabula rasa once weight
convergence reveals prior quality.

\textbf{Critical Implementation Note:}
The \texttt{selection\_token} returned by \texttt{select\_model()} must
be passed to \texttt{update()}.  Omitting it causes complete failure of
meta-learning (weights remain frozen at 50\%).

% ----------------------------------------------------------------------------
% Recommendation 2: Exploration and Gamma
% ----------------------------------------------------------------------------
\paragraph{Exploration and Mixing Defaults.}
We use constant $\alpha{=}2.0$ and $\gamma{=}0.05$ throughout.
Constant exploration preserves change-detection capability in
non-stationary environments, accepting a modest static-benchmark
premium.
The mixing parameter $\gamma$ has negligible impact on regret
(total range ${<}1.0$ across tested values) and requires no tuning.

% ----------------------------------------------------------------------------
% Summary Table Reference
% ----------------------------------------------------------------------------
Table~\ref{tab:strategy_guide} summarises our deployment recommendations.
